%!TEX root = ../../thesis.tex
\chapter{関連研究}
\label{chap:related_work}

本章では，本研究の基礎となるシナリオ駆動形ナビゲーションと，本研究で用いる基盤モデル，および言語の指示に基づくナビゲーションについて述べる．

\section{シナリオ駆動形ナビゲーション}
\label{sec:related_scenario}

シナリオ駆動形ナビゲーション\cite{scenario_si2023}は，人間が道案内に用いるような文でロボットに経路を与える手法である．
シナリオは「突き当りまで直進．左折．」のように，通路の特徴と行動を組み合わせた文の列として記述する．
ロボットはシナリオを条件と行動の組に分解し，カメラ画像から認識した通路形状が条件を満たしたときに次の行動へ移る．
これにより，目的地までの経路を，距離や座標ではなく通路の形状の並びとして与えることができる．

従来手法では，通路形状の認識と，通路に沿った走行に，視覚と行動のend-to-end学習によるモデルを用いていた\cite{scenario_si2023}．
これらのモデルは走行対象の環境で収集したデータで学習するため，外観の異なる環境で用いるには，その環境でのデータ収集と再学習が必要であった．

% TODO: 従来手法で用いたモデルの構成と，未学習環境での性能の課題を具体的に書く．

\section{視覚ナビゲーションの基盤モデル}
\label{sec:related_vnm}

視覚ナビゲーションでは，複数のロボットと環境から収集したデータを用いるViNT\cite{vint}や，Transformerと拡散モデルを用いて目標指向ナビゲーションと探索を統一的に扱う\NoMaD\cite{nomad}が提案されている．

ViNTは，現在と過去の観測画像およびゴール画像をEfficientNet\cite{efficientnet}で符号化し，Transformerで統合して，ゴールまでの時間的な距離と将来のwaypoint列を出力する\cite{vint}．
異なるロボットで収集した多様な走行データで学習することで，学習に用いていない環境やロボットにも適用できることが示されている．
ViNTでは，GPS waypointや左折・右折・直進といった高レベルの経路指示を入力とする適応も検討されているが，その指示は，事前に記録した経路の地図とその上での自己位置推定に基づいて算出される\cite{vint}．

\NoMaD は，ViNTと同様の観測エンコーダに，ゴール画像の入力を確率的に隠すゴールマスキングを導入したモデルである\cite{nomad}．
行動の生成には拡散モデル\cite{ddpm,diffusion_policy}を用い，ノイズから段階的にwaypoint列を生成することで，複数の選択肢がある場面での多峰的な行動分布を表現する．
\NoMaD はゴール画像をマスクすることでゴール非依存の探索も行えるが，その行動は特定の目的地へ向かうものではない．
本研究では，このゴールマスク時の観測表現に目標方向の情報を加えることで，ゴール画像を用いずに指定した方向へ進ませる．

\section{画像認識の基盤モデル}
\label{sec:related_vision}

画像認識においても，自己教師あり学習によって汎用的な視覚表現を獲得した\DINOvTwo \cite{dinov2}や，単眼RGB画像から絶対尺度の距離を推定する\UniDepth \cite{unidepth}が提案されている．

\DINOvTwo は，Vision Transformer（ViT）\cite{vit}を，ラベルを用いない自己蒸留によって大規模な画像データセットで学習したモデルである\cite{dinov2}．
得られた特徴は，特徴抽出器を固定したまま線形層を学習するだけで，画像分類やセグメンテーション，深度推定など多様なタスクに利用できることが示されている．
本研究では，この性質を利用し，\DINOvTwo を凍結したまま通路分類に用いる．

\UniDepth は，単眼RGB画像から，カメラの内部パラメータを与えずに絶対尺度の深度を推定するモデルである\cite{unidepth}．
カメラの情報を画像から推定するモジュールを持ち，多様なデータセットで学習することで，学習に用いていない環境でも距離を推定できる．
本研究では，\UniDepth の推定深度から，通路分類に用いる鳥瞰の占有格子と，衝突回避に用いる障害物の点群を求める．

\section{衝突回避}
\label{sec:related_care}

\CARE は，視覚ナビゲーションモデルが出力した軌道を，推定した障害物からの反発力に基づいて補正する衝突回避手法である\cite{care}．
ナビゲーションモデル自体を再学習せずに組み合わせられるため，本研究では経路追従モジュールの出力の補正に用いる．

\section{Vision-and-Language Navigation}
\label{sec:related_vln}

関連研究として，言語で記述された指示に従って移動するVision-and-Language Navigation（VLN）が挙げられる．
LM-Navは，大規模言語モデルで指示から目印を抽出し，視覚言語モデルで目印を事前に構築したトポロジカルグラフ上の画像と対応付け，視覚ナビゲーションモデルで移動する\cite{lmnav}．
NaVILAは，視覚言語モデルが「前に75\,cm進む」のような中間的な言語の行動を出力し，それを脚ロボットの移動方策が実行する\cite{navila}．
LM-Nav\cite{lmnav}やNaVILA\cite{navila}のように，言語による指示で未学習の環境を移動すること自体はすでに実現されつつある．
しかし，これらの手法は数十億以上のパラメータを持つ言語モデルや視覚言語モデルを用いており，それらの推論をロボット外部の計算機で行っている\cite{lmnav,navila}．

\section{本研究の位置付け}
\label{sec:related_position}

本研究では，シナリオの解釈をルールベースで処理し，深層学習モデルを通路分類と経路追従に限定することで，用いるモデルのパラメータ数を合計約7500万に抑え，ロボットに搭載可能な計算機での動作を想定したシステムを構成する．
また，実行時に用いるセンサは単眼RGBカメラのみである．
視覚ナビゲーション基盤モデルと比べると，ゴール画像の代わりにシナリオから決まる目標方向を与える点が異なり，走行対象の環境を事前に走行する必要がない．
VLNと比べると，自由な言語の指示は扱えないが，定型的なシナリオに限定することで，大規模な言語モデルを用いずに経路を指示できる．
